% fixture: amsart/amsart-equation-numbers
% tier: T3
% engine: document
% covers: amsart's equation numbers on the left (amsmath with leqno: \veqno is \leqno, § 1203–1205): a number beside a centred formula, a formula pushed right of a number less than twice its width away, a number too wide for the line set above the formula at the left with no \abovedisplayskip, always the long display skips (after a short line too), \label and \ref of an equation
% issue: #360
\documentclass{amsart}
\begin{document}
The breaking of paragraphs into lines is a small part of this craft, but it decides how evenly the words are spread and how easily the eye moves from one line to the next.
\begin{equation}\label{eq:pyth}
a^2+b^2=c^2
\end{equation}
A short line.
\begin{equation}
x=y
\end{equation}
A formula that comes close to the number moves to the right of the centre, as in
\begin{equation}
a_1+a_2+a_3+a_4+a_5+a_6+a_7+a_8+a_9+a_{10}+a_{11}+a_{12}+a_{13}+a_{14}+a_{15}+a_{16}+a_{17}=0
\end{equation}
and one as wide as the line leaves the number above it, as in
\begin{equation}
a_1+a_2+a_3+a_4+a_5+a_6+a_7+a_8+a_9+a_{10}+a_{11}+a_{12}+a_{13}+a_{14}+a_{15}+a_{16}+a_{17}+a_{18}+a_{19}+a_{20}=0
\end{equation}
By \eqref{eq:pyth} and (\ref{eq:pyth}) the text goes on after the displays.
\end{document}
